% Generated draft: requires mathematical and source review.
\documentclass[11pt]{article}
\usepackage[margin=25mm]{geometry}
\usepackage{fontspec}
\setmainfont{Latin Modern Roman}
\usepackage{amsmath,amssymb,mathtools,unicode-math}
\setmathfont{Latin Modern Math}
\usepackage{enumitem,needspace,xurl,hyperref}
\setlength{\parindent}{0pt}
\setlength{\parskip}{5pt}
\setlength{\emergencystretch}{3em}
\newcommand{\N}{\mathbb{N}}
\newcommand{\Z}{\mathbb{Z}}
\newcommand{\Q}{\mathbb{Q}}
\newcommand{\R}{\mathbb{R}}
\newcommand{\C}{\mathbb{C}}
\newcommand{\F}{\mathbb{F}}
\newcommand{\A}{\mathbb{A}}
\newcommand{\PP}{\mathbb{P}}
\newcommand{\E}{\mathbb{E}}
\newcommand{\eps}{\varepsilon}
\newcommand{\dd}{\,\mathrm{d}}
\begin{document}
\section*{UnsolvedMath ID: 30001054}
\textbf{Source problem code:} catalogue entry
\subsection{Definitions and mathematical statement}
Let $\Sigma$ be an allowable double-permutation sequence in the sense of the cited source. The problem asks whether there exist a family $\mathcal C$ of connected sets and a family $\mathcal A$ of pseudoline separators and common tangents whose associated sequence is $\Sigma$. In notation for the source's sequence construction, the assertion is
\[
\forall\Sigma\text{ allowable}\quad
\exists(\mathcal C,\mathcal A)\text{ satisfying the source's geometric conditions}:\qquad
\operatorname{Seq}(\mathcal C,\mathcal A)=\Sigma.
\]
Thus the target is surjectivity of the geometric encoding onto all allowable sequences, not just necessity of the allowable axioms for an existing configuration.

In a conventional planar use, a pseudoline is a simple unbounded curve playing the role of a line, with the prescribed topological crossing and separation properties. This terminology does not imply a straight line. For connected sets, ``common tangent'' must be given a supporting or topological meaning, not silently a derivative-based tangency condition. The exact ambient surface, separation and tangency rules, crossing restrictions and sweep construction are not supplied.

In particular, the record does not define double-permutations, their labels or possible sign conventions, the allowed transitions in a sequence, or the equivalence under which an encoded sequence agrees with $\Sigma$. Ordinary allowable permutation sequences cannot be substituted without justification. No convexity, compactness, pairwise disjointness or smoothness of the connected sets is imposed, since none is stated here. These missing definitions prevent an intrinsic self-contained realizability criterion, but the source's universal question and both geometric families are retained.

The supplied assessment checked 2026-08-21 labels the problem open and says the primary treatment explicitly leaves realization by connected sets with the required tangents and separators unknown. The corrected-verified note describes semantic and typographic cleanup without unsupported additions; it does not reproduce the missing definitions.
\subsection{Short English statement}
Ask whether every allowable double-permutation sequence can be realized by connected sets equipped with the required pseudoline separators and common tangents. The combinatorial axioms and geometric encoding must be supplied from the cited treatment.
\subsection{Sources}
\begin{itemize}
\item Local UnsolvedMath registry, permanent ID 30001054, source code catalogue entry; supplied complete record and corrected-verified cleanup note. Literature checked 2026-08-21: the assessment is supplied dated curation, not independently verified here.
\item Discrete Geometry (2009), Oberwolfach Reports, DOI 10.4171/owr/2008/44, pp. 2477--2556: \url{https://doi.org/10.4171/owr/2008/44}. Proposed-year metadata gives 2008.
\item Pollack: \url{https://math.nyu.edu/faculty/pollack/pubs/dblperm11-26-06.pdf}. This is the supplied primary-treatment reference, not an inferred named proposer.
\end{itemize}
\end{document}
